% Short summary coordinated with the integrated thesis, 2026-09-09.
How children develop expressive vocabulary and grammatical behavior from limited, unevenly distributed linguistic experience is a central question in cognitive science. Language models make learning assumptions experimentally tractable, but their developmental relevance depends on how experience and behavior are compared. This thesis combines a framework for calibrating child--model comparisons with four computational studies of expressive vocabulary, retrieval-supported syntax, dialogue simulation, and learning from caregiver-derived rewards. The studies use distinct learners and training regimes to examine what predictive exposure, access to stored instances, and feedback contribute to specific linguistic outcomes.

The \textsc{corpus-CDI} benchmark aligns estimated training exposure, generated production quantity, and evaluation vocabulary to compare children's expressive development with character-level LSTMs and Transformers. The tested models do not jointly reproduce children's trajectories of lexical diversity, word-form validity, and vocabulary growth; low-frequency vocabulary remains particularly difficult to produce. A separate study uses $k$-nearest-neighbor retrieval to supplement parametric language models. Retrieval improves sensitivity to syntactic test items with low mean content-word frequency and partially narrows the lexical frequency gap, with structural information playing an important role. Turning to interaction, a benchmark comparing simulated child--caregiver exchanges with naturalistic dialogue shows that selected utterance-level and age-related patterns can be approximated, while sustained conversations reveal greater divergence. Simulated caregivers show higher semantic alignment and lower semantic diversity than the human reference. In an independent reinforcement-learning study, caregiver-derived rewards have differentiated effects: structural alignment most consistently improves the grammaticality of generated utterances, without consistent improvements on BLiMP and Zorro minimal-pair tests.

Together, these findings identify selective benefits of retrieval and feedback alongside persistent mismatches across developmental trajectories and evaluation tasks. The thesis contributes production and interaction benchmarks, controlled computational interventions, and a framework connecting model outcomes to bounded developmental claims. It motivates theories that distinguish the contributions of exposure, memory, and interaction, and test their effects across multiple measures of language use.
